\ch{Chapter 4: Results and Analysis}\label{chap:Chapter_4}

This chapter presents the results of the survey data collection and analysis described in Chapter~3. The analysis executes the plan specified in the methodology: descriptive characterization, composite construction and reliability assessment, comparative testing across professional field and experience level, association analysis, and the planned bias diagnostics. The chapter organizes findings by the three research questions established in Chapter~1, answering each descriptively first and consolidating the three answers in Table~\ref{tab:answers} before it turns to the subgroup, regression, and robustness analyses that qualify them. Consistent with Chapter~3, the chapter reports parametric descriptive statistics (means, standard deviations) alongside ordinal summaries: medians and interquartile ranges accompany the item-level descriptives, medians accompany the composite descriptives, and the comparative tables report means with the rank-based tests. Non-parametric tests serve for inferential comparisons where distributional assumptions are not met. Because the sample is large, the analysis foregrounds effect sizes and confidence intervals and treats statistical significance as secondary.

A fully scripted pipeline produced all statistics, tables, and figures in this chapter from the collected response file, and independent recomputation cross-validated every reported value. The data depart from the proposal's expectations in three consequential respects: sample size, two investment-willingness items that show ceiling effects, and a non-monotonic experience pattern. The chapter reports these departures directly because they bear on how the findings should be interpreted.

\section{Data Preparation and Sample}

Data collection yielded \textbf{1{,}600} submitted responses, of which \textbf{1{,}411} remained for analysis. A pre-registered quality screen removed 188 abandoned or time-aberrant submissions: those answering at most the two opening questions, completing in under three minutes, or leaving the survey open beyond five days. A post-hoc response-authenticity screen (Section~\ref{sec:forensic-validation}) then removed one further submission from a coordinated automated cluster. Every count and statistic in this chapter refers to that analytic sample of 1{,}411, except where item-level missingness reduces the available $n$, as noted. The pipeline decoded Likert items from their anchored string form (for example, ``4 - Agree'') to the integer 1--5 scale and retained the binary and ternary items as categorical factors. Within the analytic sample, complete data across all twenty-three Likert-scaled items are present for 1{,}390 respondents (98.5\%). Available-case analysis handled item-level missingness, which never exceeds six cases on any single item, consistent with the Chapter~3 plan to assess rather than impute missing data. Because the removed submissions answered at most the two opening questions---most of them nothing at all---and clustered at implausible completion times, excluding them sharpens rather than biases the analytic sample. It strips abandonment noise that would otherwise distort the composite distributions and the completion-time-based wave diagnostic, without privileging any particular pattern of substantive response.

At 1{,}411 completed responses, the realized sample is roughly three times the planned target of 385--450 derived in Chapter~3. This far exceeds the requirement for a five-percent margin of error and strengthens the precision of population estimates. It also reframes the inferential reading anticipated in the methodology: at $n>1{,}400$, virtually any non-trivial difference attains $p<.001$, so statistical significance carries little interpretive weight on its own. What the sample delivers is precision. The 95\% margin of error on the headline proportions is roughly one to two percentage points, for example $\pm1.2$ points on the 93.9\% Section~2 endorsement and $\pm2.2$ points on the 78.2\% low-familiarity share. The professional-field comparison is likewise powered essentially at unity. With 501 respondents in each focal group, the design detects an effect as small as $d=0.2$ with power near $0.89$, rising above $0.99$ for any effect of $d\geq0.27$, at $\alpha=.05$, far below the smallest effect observed. The analysis below therefore foregrounds effect sizes (Cohen's $d$, the rank-biserial correlation $r_{rb}$, Cram\'er's $V$, and $\eta^2$) and treats $p$-values as secondary. Because at this sample size every non-trivial test returns a $p$-value below any conventional threshold, the chapter reports significant results as $p<.001$ and gives exact $p$-values for non-significant ones. The effect sizes and their confidence intervals carry the inferential weight.

Two demographic features of the realized sample differ from the proposal's framing. First, the professional-field item (Q6.3) resolved into \emph{four} categories rather than the planned Information Technology versus Information Assurance dichotomy: Information Technology (35.5\%), Information Assurance/Cyber Security (35.5\%), Engineering (24.0\%), and Other (5.0\%). Following the Chapter~3 plan, the primary field comparison below contrasts Information Technology against Information Assurance/Cyber Security. The Engineering and Other groups are retained for descriptive context but not forced into a binary that the instrument did not produce. Second, the experience distribution is weighted toward earlier-career respondents (40.1\% at 1--5 years; 34.7\% at 6--10 years; 24.0\% at 11--15 years), and the 16+ year cohort contains only 18 respondents (1.3\%). These percentages are computed on the full analytic sample of 1{,}411 respondents. As Chapter~3 foresaw, that cell is too small to support stable pairwise inference, and the results below bear out this limitation.

\section{Research Question 1: Awareness of Digital Engineering}

Research Question~1 addresses awareness. The analysis draws on the Section~1 awareness items and the Section~2 capability-recognition items, following the layered assessment in Chapter~3 that separates familiarity with the named discipline from recognition of its capabilities once those capabilities are described in functional terms.

\subsection{Self-Reported Familiarity and Exposure}

Self-reported familiarity with Digital Engineering (Q1.1) is low across the sample. Nearly half of respondents (49.6\%) selected ``Not familiar'' and a further 28.6\% ``Slightly familiar,'' placing 78.2\% in the bottom two categories; the item mean is 1.94 ($SD=1.20$, $Mdn=2$). Only 17.6\% (95\% CI [15.7, 19.6]) rated themselves ``Very'' or ``Extremely'' familiar (Figure~\ref{fig:fam}). Recent professional exposure (Q1.2) is consistent: 22.2\% reported encountering Digital Engineering methodologies, frameworks, or tools in their work within the past two years, against 77.8\% who reported none. The familiarity self-assessment thus answers the direct, surface-level component of Research Question~1 unambiguously: awareness of Digital Engineering, as a named discipline, is low within this population of IT and Information Assurance professionals.

Beyond its central tendency, the shape of that distribution carries information of its own. Familiarity is bimodal rather than uniformly low. Two lowest categories hold 78.2\% of respondents and the middle category holds only 4.2\%, yet ``Very familiar'' recovers to 13.6\%, more than three times the middle category. Responses therefore separate into a large stratum with little or no acquaintance with the discipline and a much smaller stratum that knows it well, with few respondents occupying the ground between them. That structure bears on the comparative analyses, because a sample composed of two such strata can produce group differences that reflect the composition of each group as much as the professional field defining it. The robustness analyses later in this chapter test that possibility directly and find that it accounts for part of the field difference while leaving most of it intact.

\begin{figure}[H]
  \centering
    \includegraphics[width=0.78\linewidth, alt={Bar chart of self-reported familiarity with Digital Engineering (Q1.1) on a five-point scale. The response mass concentrates in the two lowest categories, Not familiar and Slightly familiar, which together hold 78.2 percent of respondents; the item mean is 1.94.}]{figures/fig1_familiarity.pdf}
  \caption{Self-reported familiarity with Digital Engineering (Q1.1).}
  \label{fig:fam}
\end{figure}

\subsection{Recognition of Specific Capabilities}

The Section~2 capability-recognition items contrast sharply with the familiarity data. Agreement with the six capability statements is uniformly high: 93.9\% of Section~2 responses fall at ``Agree'' or ``Strongly agree,'' with item means between 4.29 and 4.50 (Table~\ref{tab:desc}, items Q2.1--Q2.6). The same respondents who reported little familiarity with Digital Engineering nonetheless affirmed, at very high rates, that Digital Engineering ``includes model-based systems engineering approaches'' (Q2.1, 95.7\% agreement), ``can enable digital twin development'' for IT systems (Q2.2, 95.8\%), and ``can improve security control implementation through automated compliance checking'' (Q2.6, 95.7\%). The full distribution of the twenty-two agreement-scale items (Sections~2 through~5), centered on the neutral midpoint, appears in Figure~\ref{fig:div}: every item is strongly right-shifted, with minimal disagreement throughout. Because Q1.1 measures familiarity on a different response scale, its distribution appears separately in Figure~\ref{fig:fam} rather than combined with the agreement items.

\begin{figure}[H]
  \centering
    \includegraphics[width=0.96\linewidth, alt={Diverging bar chart of the twenty-two agreement-scale items from Sections 2 through 5, centered on the neutral midpoint. Every item is strongly right-shifted toward agreement, with minimal disagreement throughout; Section 2 agreement reaches 93.9 percent with item means between 4.29 and 4.50. The familiarity item (Q1.1) uses a different response scale and is shown separately in the familiarity figure.}]{figures/fig2_diverging.pdf}
  \caption{Response distributions for the twenty-two agreement-scale items (Q2.1--Q5.6).}
  \label{fig:div}
\end{figure}

This divergence is the central finding for Research Question~1, and it bears out the distinction drawn in Chapter~3 between awareness of the named discipline (Section~1) and recognition of described capabilities (Section~2). Two readings are admissible, and the data cannot fully separate them. The charitable reading is that respondents recognize and endorse the described capability relationships once those relationships are stated in plain functional language, even without prior familiarity with the umbrella term. The cautionary reading is that the Section~2 stems describe capabilities in affirmative terms that invite agreement, so high endorsement may partly reflect acquiescence or face-validity assent rather than substantive knowledge. Neither reading can be settled here, because the instrument carries no knowledge-verification item against which an agreement response could be scored correct or incorrect. Section~2 therefore bounds what Research Question~1 can claim: it measures recognition, and comprehension lies outside its reach. The near-absence of disagreement on these items, visible in Figure~\ref{fig:div}, is consistent with either account and should temper any claim that Section~2 shows pre-existing domain awareness.

A straight-lining check across the twenty-two agreement-scale items quantifies the acquiescence risk: 1.6\% of respondents (22 of the 1{,}390 with the complete battery) gave an identical response to all twenty-two items, a small yea-saying signature (Appendix~B). Most vary their responses (median within-respondent standard deviation 0.50), so crude and uniform acquiescence is a minor presence rather than a driver of the high endorsement. That check bounds the extreme form and not acquiescence as such, because a mild yea-saying tilt would produce exactly the within-respondent variation observed here, and Chapter~5 accordingly keeps it among the open accounts. The present results therefore characterize Section~2 as measured \emph{endorsement of capability descriptions} rather than as evidence of prior awareness. Taken with the Section~1 findings, the answer to Research Question~1 is that awareness of Digital Engineering as a discipline is low, while endorsement of its constituent capabilities, once described, is high.

\begin{table}[ht]
\centering
\caption{Descriptive statistics for the Likert-scale items (1--5).}
% ALT: Descriptive statistics for Likert-scale items (1--5). Agreement = score $\geq$ 4; intervals are 95\% Wilson bounds.
\label{tab:desc}
\small
\setlength{\tabcolsep}{4pt}
\begin{tabular}{@{}llrrrccrc@{}}
\toprule
Item & Construct & $n$ & $M$ & $SD$ & Mdn & IQR & Agree\% & 95\% CI \\
\midrule
Q1.1 & Familiarity with DE (Q1.1) & 1{,}411 & 1.94 & 1.20 & 2 & 1--2 & 17.6 & [15.7, 19.6] \\
Q2.1 & MBSE improves dev processes & 1{,}411 & 4.42 & 0.57 & 4 & 4--5 & 95.7 & [94.6, 96.7] \\
Q2.2 & Digital twin for IT systems & 1{,}405 & 4.38 & 0.56 & 4 & 4--5 & 95.8 & [94.6, 96.7] \\
Q2.3 & Continuous integration / data-driven & 1{,}407 & 4.39 & 0.57 & 4 & 4--5 & 95.8 & [94.6, 96.7] \\
Q2.4 & Digital twin for security scenarios & 1{,}407 & 4.50 & 0.58 & 5 & 4--5 & 95.7 & [94.5, 96.7] \\
Q2.5 & Continuous security validation & 1{,}408 & 4.29 & 0.72 & 4 & 4--5 & 84.6 & [82.6, 86.4] \\
Q2.6 & Automated compliance checking & 1{,}407 & 4.41 & 0.57 & 4 & 4--5 & 95.7 & [94.5, 96.7] \\
Q3.1 & Relevant to IT sector & 1{,}410 & 4.27 & 0.72 & 4 & 4--5 & 84.6 & [82.6, 86.4] \\
Q3.2 & Relevant to IA challenges & 1{,}408 & 4.39 & 0.57 & 4 & 4--5 & 95.7 & [94.6, 96.7] \\
Q3.3 & Digital twins vs.\ production & 1{,}407 & 4.51 & 0.58 & 5 & 4--5 & 95.8 & [94.6, 96.7] \\
Q3.4 & Models map environment/config & 1{,}407 & 4.42 & 0.57 & 4 & 4--5 & 95.9 & [94.7, 96.8] \\
Q3.5 & Lifecycle mgmt for compliance & 1{,}408 & 4.30 & 0.72 & 4 & 4--5 & 84.8 & [82.8, 86.6] \\
Q3.6 & Org has compliance needs & 1{,}408 & 3.89 & 0.81 & 4 & 3--5 & 61.1 & [58.5, 63.6] \\
Q4.1 & Value to IT processes & 1{,}411 & 4.39 & 0.57 & 4 & 4--5 & 95.9 & [94.7, 96.8] \\
Q4.2 & Reduce dev cycle time & 1{,}407 & 4.29 & 0.70 & 4 & 4--5 & 85.5 & [83.6, 87.2] \\
Q4.3 & Improve quality / fewer defects & 1{,}408 & 4.30 & 0.72 & 4 & 4--5 & 84.8 & [82.8, 86.6] \\
Q4.4 & Improve dev collaboration & 1{,}408 & 4.12 & 0.64 & 4 & 4--5 & 84.7 & [82.7, 86.4] \\
Q5.1 & Value to IA/cyber operations & 1{,}409 & 4.12 & 0.64 & 4 & 4--5 & 85.0 & [83.0, 86.7] \\
Q5.2 & Faster vuln.\ remediation & 1{,}408 & 4.13 & 0.80 & 4 & 3--5 & 73.8 & [71.4, 76.0] \\
Q5.3 & Improve security posture & 1{,}408 & 4.31 & 0.72 & 4 & 4--5 & 84.8 & [82.8, 86.6] \\
Q5.4 & Enhance threat modeling & 1{,}408 & 4.32 & 0.75 & 4 & 4--5 & 83.0 & [81.0, 84.9] \\
Q5.5 & Cross-team security collab. & 1{,}408 & 4.22 & 0.71 & 4 & 4--5 & 83.3 & [81.3, 85.2] \\
Q5.6 & Better framework compliance & 1{,}408 & 4.23 & 0.71 & 4 & 4--5 & 83.5 & [81.5, 85.4] \\
\bottomrule
\end{tabular}
\end{table}


\section{Research Question 2: Perceived Value}

Perceived value is high across all three value blocks, which is the answer the second research question receives. Three groups of items measure it: applicability in Section~3, information technology value in Section~4, and security value in Section~5. Agreement runs at 86.3\% for applicability, 87.7\% for information technology, and 82.2\% for security, and item means fall between 3.9 and 4.5 (Table~\ref{tab:desc}). Item-level agreement with 95\% Wilson confidence intervals~\parencite{Wilson_1927} appears in Figure~\ref{fig:agree}. The intervals are narrow throughout, a consequence of the large sample, though that narrowness is conditional on the accessible sample rather than on a random draw of the profession.

A conservative reading restricts attention to the low-familiarity majority, for whom the value finding cannot be attributed to prior enthusiasm. Among the 78.2\% of respondents in the two lowest familiarity categories (Q1.1 $\in \{1,2\}$), perceived value remains high, with composite means of 4.21 for IT Value and 4.08 for Security Value, only modestly below the whole-sample means of 4.33 and 4.20. The central finding, that value is perceived as high even where familiarity is absent, therefore does not rest on the more familiar or more favorably disposed respondents. The same caution that qualifies the Section~2 items applies here: because every value item is affirmatively keyed, high agreement within the low-familiarity stratum rules out prior enthusiasm but cannot separate genuine appraisal from face-validity assent by respondents new to the concept. The value finding is therefore best read as receptiveness consistent with genuine appraisal rather than as demonstrated conviction.

\begin{figure}[tbp]
  \centering
  \includegraphics[width=0.96\linewidth, alt={Bar chart of item-level agreement, the percentage of respondents scoring four or higher, for each survey item grouped by section, shown with 95 percent Wilson confidence interval whiskers. Agreement is high across all sections and the intervals are narrow because of the large sample.}]{figures/fig4_agreement.pdf}
  \caption{Item-level agreement (score $\geq 4$) by section with 95\% Wilson confidence intervals.}
  \label{fig:agree}
\end{figure}

The two investment-willingness items require careful handling. IT investment willingness (Q4.5) records 1{,}192 ``Yes'' (84.5\%), 218 ``Unsure'' (15.5\%), and a single ``No'' ($<$0.1\%) across the sample. Security investment willingness (Q5.7) records 99.8\% ``Yes'' (1{,}408 of 1{,}411), with one ``No'' and two ``Unsure.'' These distributions are at or near the response ceiling. Chapter~3 intended these items as evidence of organizational receptiveness, and at face value they supply it overwhelmingly. A near-unanimous item, however, has almost no variance, so it cannot discriminate among respondents or support meaningful subgroup contrasts. Such uniformity is also a recognized signature of response-set behavior or of a sample skewed toward the favorably disposed. Q5.7 therefore establishes a near-universal stated willingness \emph{conditional on demonstrated return on investment} while contributing essentially nothing to differentiation among respondents. This pattern should be read as conviction rather than discriminating evidence. Interest and intent items reinforce the value findings: 89.0\% of respondents indicated interest in learning more about Digital Engineering applications (Q6.1), and 95.9\% would recommend that their organization explore Digital Engineering adoption (Q6.2).
% Verified: Q6.1 = 89.0% (n = 1,403) and Q6.2 = 95.9% (n = 1,397) reproduced in Analysis/DigitalEngineering_Analysis.ipynb (Section 3, interest/recommendation cell).
Research Question~2 therefore receives a precisely estimated answer: perceived value is high, and professionals rate Digital Engineering capabilities as valuable across the applicability, IT, and security domains, with the qualification that the investment-willingness items reflect near-unanimous conviction rather than graduated discrimination. Every statistic in this section is a perception. They establish that this population expects Digital Engineering to be valuable in IT and security work, and whether Digital Engineering delivers that value once implemented is an implementation question this design does not address.

\section{Research Question 3: Anticipated Benefit for Job Performance, Compliance, and Capability}

Research Question~3 asks to what extent professionals believe Digital Engineering could help them perform their jobs, meet compliance requirements, or enhance organizational capabilities. The analysis draws on the operational and compliance-oriented items across Sections~3--5.

Research Questions~2 and~3 draw on overlapping survey material, and saying so plainly matters for how their two answers are counted. Chapter~1 distinguishes them conceptually, with Research Question~2 addressing general perceived value and Research Question~3 addressing applicability to the respondent's own job performance, compliance obligations, and organizational capability. The instrument operationalizes both through the Section~3 to Section~5 items, regrouped rather than separately measured, so the two answers rest on a common evidentiary base and are best read as one appraisal examined at two levels of specificity.

Agreement is substantial on all three groupings. The IT operational items, Q4.1 through Q4.4, cover meaningful value, cycle time, quality, and collaboration, and they draw 87.7\% agreement. The security operational items, Q5.1 through Q5.5, cover value to operations, vulnerability remediation, posture, threat modeling, and cross-team collaboration, and they draw 82.0\%. The three explicitly compliance-framed items, Q3.5, Q3.6, and Q5.6, draw 76.5\%.

Those three aggregates conceal the finding that matters most for Research Question~3, because each averages over items that behave differently. Disaggregating the compliance grouping makes the point plainly. Its two capability-framed items draw agreement in line with the rest of the instrument: Q3.5, on lifecycle management for compliance activity, reaches 84.8\% ($M=4.30$), and Q5.6, on improved framework compliance, reaches 83.5\% ($M=4.23$). Q3.5 pairs compliance with service delivery in a single stem, so some of its agreement may attach to the second clause, which makes Q5.6 the cleaner of the two comparators. Its third item behaves differently. Q3.6 (``my organization faces regulatory or compliance requirements that could benefit'') draws 61.1\% agreement ($M=3.89$; Table~\ref{tab:desc}), more than twenty percentage points below its companions. The 76.5\% aggregate is therefore an average across a genuine split rather than a uniform level of reservation, and reporting it alone would obscure what the items separately show.

What separates Q3.6 from its companions is the kind of judgment it requires. Q3.5 and Q5.6 ask respondents to appraise what Digital Engineering can do; Q3.6 asks them to characterize their own organization's regulatory exposure. The same pattern recurs in the security grouping, where Q5.2 (faster vulnerability remediation) draws 73.8\% ($M=4.13$) against 83.0\% to 85.0\% for the other four Section~5 operational items. Q5.2 likewise asks for a claim about a measurable operational outcome rather than about a capability in principle. Its stem, however, bundles two operations, identification and remediation, so a respondent who credits Digital Engineering with the first but not the second would also hedge. The design cannot separate that reading from the own-circumstances reading, and Q5.2 is reported here as consistent with the pattern rather than as independent support for it.

Dispersion statistics corroborate the reading. Across the twenty-two agreement-scale items, Q3.6 and Q5.2 are the two lowest in agreement, and they are also the only two whose interquartile range extends down to the neutral midpoint (3--5 against 4--5 for every other item; Table~\ref{tab:desc}). The next lowest item, Q5.4, sits at 83.0\%, more than nine points above Q5.2. Reservation in this instrument is therefore concentrated in two items rather than distributed across the compliance and security blocks. Q3.6 carries the clearer case, because its stem is single-barreled and asks only for a characterization of the respondent's own organization. Read together, the two items locate reservation in claims respondents must ground in their own circumstances, with the weight of that reading resting on Q3.6.

Read together, the Research Question~3 items support the proposition that respondents believe Digital Engineering could help in their work, subject to one qualification the disaggregation makes precise. Endorsement is near-uniform wherever the instrument asks what the capability could do, and it falls by ten to twenty points wherever the instrument asks respondents to assert something concrete about their own organization or its measurable performance. This qualification bears on how the answer should be used. Anticipated benefit is high and is evidence of receptiveness; it is a weaker basis for inferring that respondents see the described benefits as already applicable to their own circumstances, which is the belief that would most directly precede adoption.

\section{Direct Answers to the Research Questions}

All three research questions ask to what extent, so each admits a descriptive answer, and Table~\ref{tab:answers} states the three before the analysis turns to how they are distributed. Awareness of Digital Engineering as a named discipline is low: 78.2\% of respondents fall in the two lowest familiarity categories and 22.2\% report recent professional exposure. Perceived value is high and precisely estimated across all three value blocks, at 86.3\%, 87.7\%, and 82.2\% agreement. Anticipated benefit for job performance, compliance, and organizational capability is likewise high, at 87.7\% on the IT operational items and 82.0\% on the security operational items, and it is most reserved on the explicitly compliance-framed items at 76.5\%. These are the answers the study returns.

\begin{table}[htbp]
  \centering
  \caption{Direct answers to the three research questions.}\label{tab:answers}
  % ALT: Table stating the descriptive answer to each of the three research questions alongside the key evidence supporting it.
  \small
  \begin{tabular}{@{}>{\raggedright\arraybackslash}p{0.19\textwidth}>{\raggedright\arraybackslash}p{0.28\textwidth}>{\raggedright\arraybackslash}p{0.45\textwidth}@{}}
    \toprule
    \textbf{Research question} & \textbf{Direct answer} & \textbf{Key evidence} \\
    \midrule
    RQ1 Awareness & Low awareness, near-universal capability endorsement & Familiarity mean 1.94; 78.2\% in the two lowest categories; 22.2\% recent exposure; 93.9\% Section~2 agreement \\
    \addlinespace
    RQ2 Perceived value & High perceived value across all value blocks & Agreement 86.3\% (applicability), 87.7\% (IT), 82.2\% (security); composite means 4.0--4.3 \\
    \addlinespace
    RQ3 Anticipated benefit & High anticipated benefit, most reserved on organization-specific compliance & IT operations 87.7\%, security operations 82.0\%, compliance-framed 76.5\% \\
    \bottomrule
  \end{tabular}
  \smallskip

  {\footnotesize \textit{Note.} All figures are from the authenticated analytic sample ($N=1{,}411$). Each answer is descriptive. The subgroup, regression, and robustness analyses later in this chapter qualify these answers without displacing them.}
\end{table}

The remainder of the chapter examines how those answers are distributed and how firmly they hold. Composite, comparative, regression, and robustness analyses characterize variation within the population and test the findings against alternative explanations. They refine the three answers above rather than displace them.

\section{Composite Scores and Reliability}

Item averaging produced the three composite scores specified in Chapter~3: an Awareness composite (Q1.1, Q2.1--Q2.6), an IT Value composite (Q3.1, Q3.3--Q3.5, Q4.1--Q4.4), and a Security Value composite (Q3.2, Q3.6, Q5.1--Q5.6). All three comfortably exceed the conventional Cronbach's $\alpha \geq 0.70$ threshold (Table~\ref{tab:comp}): Awareness $\alpha=0.792$, IT Value $\alpha=0.858$, and Security Value $\alpha=0.858$. Internal consistency is therefore adequate for the comparative analyses that follow.

One conceptual note bears on the Awareness composite. It averages the low-scoring familiarity item (Q1.1, $M=1.94$) with the high-scoring Section~2 capability-recognition items ($M\approx4.4$), producing a composite mean of 4.04 pulled upward by those items. The composite is internally consistent, but its central tendency should not be read as evidence that respondents are highly aware: it blends low familiarity with high capability-endorsement, and its two constituents point in opposite directions. Reporting Q1.1 separately alongside the composite, as Table~\ref{tab:desc} does, prevents that misreading.

The item statistics quantify how far apart those constituents sit. Q1.1 correlates with the mean of the six capability-recognition items at $r=0.32$, and its corrected item-total correlation is 0.315 against a range of 0.53 to 0.67 for the six Section~2 items. Deleting Q1.1 would raise the composite's $\alpha$ from 0.792 to 0.863. Under the conventional item-retention criteria, which flag items whose corrected item-total correlation falls below roughly 0.40 and whose removal improves $\alpha$, Q1.1 is therefore a deletion candidate, and the remaining six items would form a cleaner scale measuring capability endorsement. The analysis retains all seven items regardless, because Chapter~3 specified this composite before data collection and revising a pre-specified measure after seeing the results would trade the discipline of pre-specification for analytic flexibility. That decision carries one consequence the comparative sections observe: the composite indexes capability endorsement more heavily than familiarity, so the field difference reported for Awareness below is better understood as a difference in capability endorsement, while Q1.1 carries the familiarity finding on its own.

\begin{table}[ht]
\centering
\caption{Composite scores and internal consistency.}
% ALT: Composite scores and internal consistency. Cronbach's $\alpha \geq 0.70$ is the conventional acceptability threshold.
\label{tab:comp}
\small
\begin{tabular}{lrrrrcc}
\toprule
Composite & Items & $M$ & $SD$ & Mdn & 95\% CI & $\alpha$ \\
\midrule
Awareness (Q1.1, Q2.1--2.6) & 7 & 4.044 & 0.478 & 4.000 & [4.019, 4.069] & 0.792 \\
IT Value (Q3.1,3.3--3.5; Q4.1--4.4) & 8 & 4.325 & 0.463 & 4.375 & [4.301, 4.349] & 0.858 \\
Security Value (Q3.2,3.6; Q5.1--5.6) & 8 & 4.202 & 0.509 & 4.250 & [4.176, 4.229] & 0.858 \\
\bottomrule
\end{tabular}
\end{table}


\section{Comparative Analysis}

Comparative analyses examine whether awareness levels and value perceptions differ across professional field and experience level. Shapiro--Wilk tests rejected normality for all three composites (see Appendix~B), an expected result at this sample size, where the test is highly powered against trivial departures; the composite distributions are also left-skewed. Mann--Whitney $U$ and Kruskal--Wallis tests therefore serve as the primary inferential procedures, with parametric statistics reported for continuity with the technology-acceptance literature.

\subsection{Professional Field: Information Technology versus Information Assurance}

Following the analytical plan, the analysis compared Information Technology ($n=501$) and Information Assurance/Cyber Security ($n=501$) respondents on each composite. The differences are large and consistent in direction: IT respondents score substantially higher than IA/Cyber respondents on all three composites (Table~\ref{tab:field}), with Cohen's $d$ between 1.64 and 1.94 and rank-biserial correlations of magnitude $0.74$--$0.82$. The 95\% confidence intervals on $d$ are narrow and lie entirely above 1.4 (for example, Awareness $d=1.64$, 95\% normal-theory CI $[1.50,1.78]$; the bootstrap intervals in the robustness analyses below agree to within a hundredth), so the effect is both very large by conventional benchmarks and precisely estimated.

Three framings of that magnitude deserve to be kept apart. The rank-biserial coefficients carry the negative sign of the $U$-based convention while the tables report Cohen's $d$ as positive; both indicate Information Technology scoring higher. Because the higher-scoring Information Technology group sits near the scale ceiling, its variance is compressed, so Cohen's $d$, which standardizes by the pooled standard deviation, is upward-biased here. The distribution-free rank-biserial correlations are the more trustworthy effect sizes, reported alongside $d$ throughout, and the headline magnitude should be read accordingly. Expressing the same differences in the units respondents actually used supplies a third and more interpretable reading: the Information Technology advantage is 0.65 scale points on Awareness, 0.69 on IT Value, and 0.75 on Security Value, or 16 to 19 percent of the four-point interval the scale spans. The standardized and raw framings describe one gap, and reporting both keeps the standardized magnitude from carrying more rhetorical weight than the underlying difference supports.

A further caution applies to reading the two value effect sizes as independent corroboration. The construct-validity analysis below finds the IT Value and Security Value composites not discriminant from each other (heterotrait--monotrait ratio of 0.99), so their two effect sizes, identical at $d=1.94$, are better read as one result expressed on two overlapping item sets than as two confirmations of it.

The direction of the difference merits emphasis because it runs against the grain of the dissertation's framing. The Information Assurance audience the research most seeks to reach reports markedly \emph{lower} awareness and perceived value than the IT audience. The familiarity gap is the largest single contrast in the data: IT respondents average 2.88 on Q1.1 against 1.20 for IA/Cyber, and the gap propagates through the value composites (Figure~\ref{fig:box}).

An effect of this magnitude warrants scrutiny in its own right, because separations approaching two standard deviations between two adjacent professional communities are uncommon and invite an artifactual reading. Five candidate explanations warrant consideration. The first is a genuine awareness and exposure deficit within the Information Assurance community. The second is a recruitment-frame effect, in which different channels reached the two groups. The third is differential item functioning, in which security professionals calibrate affirmatively framed capability statements more conservatively. The fourth is self-categorization, in which the ``Information Assurance/Cyber Security'' label attracts more skeptical respondents. The fifth is calibrated domain skepticism, in which Information Assurance respondents discount affirmatively framed claims about security capability because their expertise equips them to appraise such claims critically. The last of these differs from the others in kind rather than in degree, because it treats the lower Information Assurance ratings as informed judgment rather than as measurement artifact or absent knowledge. One observation in the data is consistent with it: Information Assurance respondents rate the security value of Digital Engineering (3.82) below the rating Information Technology respondents give that same security value (4.57), so on the subject matter closest to their own practice the specialists are the more reserved group. That observation is consistent with calibrated skepticism without being distinctive to it, because every account listed above predicts lower Information Assurance ratings. It narrows nothing on its own and is reported as a pointer for the replication the limitations call for.

As the diagnostics below and Table~\ref{tab:threats} show, these accounts do not stand on equal footing. The within-cohort and regression analyses below exclude the experience confound, the invariance screen and the rank-based effect sizes weaken the differential-item-functioning and ceiling-artifact accounts, and the recruitment-frame and self-categorization accounts remain open because the anonymous instrument recorded no variable that could test them. The data cannot adjudicate among those open accounts; Chapter~5 weighs them. This chapter foregrounds the finding rather than burying it in Chapter~5, because it directly motivates the targeted professional-development argument the dissertation advances while remaining explicitly open about its possible sources. Table~\ref{tab:threats} consolidates the candidate explanations, the diagnostic applied to each, and what the data can and cannot exclude, so that readers can check at a glance the basis for reading the field gap as strong but provisional.

\begin{table}[htbp]
  \centering
  \caption{Candidate explanations for the Information Technology versus Information Assurance field gap, the diagnostic applied to each, and what the data can and cannot exclude.}
  % ALT: Table of candidate explanations for the IT-versus-IA field gap; each row gives the explanation, the diagnostic applied, and whether the data exclude it.
  \label{tab:threats}
  \small
  \begin{tabular}{@{}p{0.26\linewidth}p{0.34\linewidth}p{0.30\linewidth}@{}}
    \toprule
    \textbf{Candidate explanation} & \textbf{Diagnostic applied} & \textbf{Outcome} \\
    \midrule
    Experience-imbalance confound & Within-cohort models; field-by-experience association ($V=0.05$) & \textbf{Excluded}: IT exceeds IA within every experience cohort; field and experience are only weakly associated \\
    \addlinespace
    Differential item functioning & Zumbo linear DIF screen on all 23 items & \textbf{Weakened}: 21 of 23 items show no DIF; residual item differences are small relative to the composite gap \\
    \addlinespace
    Ceiling-induced variance compression inflating $d$ & Distribution-free rank-biserial correlation reported alongside $d$ & \textbf{Weakened}: the gap is large on the distribution-free measure as well ($|r_{rb}|$ of 0.74--0.82) \\
    \addlinespace
    Prior-exposure imbalance & Recent-exposure indicator added to the value regressions & \textbf{Weakened}: the field coefficient is essentially unchanged \\
    \addlinespace
    Familiarity-stratum composition & High-familiarity stratum indicator added to the value regressions & \textbf{Weakened}: the field coefficient retains 87 percent of its magnitude; 86 percent with exposure also held constant \\
    \addlinespace
    Calibrated domain skepticism & Not separable from an awareness deficit in an anonymous cross-section, since both predict lower Information Assurance endorsement & \textbf{Not excluded}: separating informed discounting from unfamiliarity requires knowledge-verification items the instrument did not carry \\
    \addlinespace
    Recruitment-frame difference & No employment-sector or recruitment-channel variable was recorded & \textbf{Not excluded}: untestable from these data; bounded by argument rather than measured \\
    \addlinespace
    Self-categorization or dispositional difference by label & Not directly testable in an anonymous cross-section & \textbf{Not excluded}: requires a design that separates label choice from field \\
    \bottomrule
  \end{tabular}
\end{table}

\begin{table}[ht]
\centering
\caption{Field comparison of the value composites, Information Technology vs.\ Information Assurance/Cyber Security.}
% ALT: Professional field comparison, Information Technology ($n=501$) vs.\ Information Assurance/Cyber Security ($n=501$). Mann--Whitney $U$ two-sided; $d$ = Cohen's $d$ with 95\% confidence interval (Hedges--Olkin large-sample approximation); $r_{rb}$ = rank-biserial correlation.
\label{tab:field}
\small
\begin{tabular}{lrrrrrcr}
\toprule
Composite & IT $M$ & IA $M$ & $U$ & $p$ & $d$ & 95\% CI $d$ & $r_{rb}$ \\
\midrule
Awareness & 4.379 & 3.732 & 218{,}766 & $<.001$ & 1.64 & [1.50, 1.78] & -0.743 \\
IT Value & 4.669 & 3.980 & 227{,}984 & $<.001$ & 1.94 & [1.79, 2.09] & -0.817 \\
Security Value & 4.566 & 3.820 & 227{,}814 & $<.001$ & 1.94 & [1.79, 2.09] & -0.815 \\
\bottomrule
\end{tabular}
\end{table}


\begin{figure}[H]
  \centering
    \includegraphics[width=0.95\linewidth, alt={Box plots of the three composite scores (Awareness, IT Value, Security Value) by professional field, with diamonds marking group means. The Information Technology and Information Assurance groups separate visibly across all three constructs, Information Technology scoring higher, with Cohen's d between 1.64 and 1.94.}]{figures/fig3_field_box.pdf}
  \caption{Composite score distributions by professional field.}
  \label{fig:box}
\end{figure}

At the item level, the analysis applied the Holm--Bonferroni family-wise correction specified in Chapter~3 within the Awareness family (Section~1 and Section~2 items). Every item survives correction with Information Technology exceeding Information Assurance throughout (Table~\ref{tab:holm_aware}); the correction overturns no conclusion, because the unadjusted $p$-values are already extreme, the expected behavior of Holm--Bonferroni when effects are large and uniform. The same procedure applied to the applicability, IT-value, and security-value families yields the same pattern: every item shows Information Technology exceeding Information Assurance and survives Holm correction (Tables~\ref{tab:holm_applic}, \ref{tab:holm_itval}, and~\ref{tab:holm_secval}). The separation is broadly uniform across the three families (mean $|r_{rb}|$ of approximately 0.56 on the applicability items, 0.56 on the security-value items, and 0.53 on the IT-value items) and slightly smaller on the IT-value items, where Q4.4 shows the smallest effect in that family ($r_{rb}=-0.45$). Together with the organization-specific compliance item Q3.6 ($r_{rb}=-0.45$), these are the smallest single-item contrasts across the three value families.

\begin{table}[ht]
\centering
\caption{Item-level field comparison within the Awareness family (IT vs.\ IA).}
% ALT: Item-level field comparison within the Awareness family (IT vs.\ IA), Mann--Whitney $U$ with Holm--Bonferroni family-wise correction ($\alpha=0.05$).
\label{tab:holm_aware}
\small
\begin{tabular}{lrrrrc}
\toprule
Item & IT $M$ & IA $M$ & $r_{rb}$ & $p_{\mathrm{Holm}}$ & Sig. \\
\midrule
Q1.1 & 2.88 & 1.20 & -0.775 & $<.001$ & \checkmark \\
Q2.1 & 4.65 & 4.17 & -0.415 & $<.001$ & \checkmark \\
Q2.2 & 4.60 & 4.15 & -0.395 & $<.001$ & \checkmark \\
Q2.3 & 4.61 & 4.16 & -0.398 & $<.001$ & \checkmark \\
Q2.4 & 4.73 & 4.25 & -0.416 & $<.001$ & \checkmark \\
Q2.5 & 4.53 & 4.02 & -0.365 & $<.001$ & \checkmark \\
Q2.6 & 4.65 & 4.17 & -0.428 & $<.001$ & \checkmark \\
\bottomrule
\end{tabular}
\end{table}

\begin{table}[ht]
\centering
\caption{Item-level field comparison within the Applicability family (IT vs.\ IA).}
% ALT: Item-level field comparison within the Applicability family (IT vs.\ IA), Mann--Whitney $U$ with Holm--Bonferroni family-wise correction ($\alpha=0.05$).
\label{tab:holm_applic}
\small
\begin{tabular}{lrrrrc}
\toprule
Item & IT $M$ & IA $M$ & $r_{rb}$ & $p_{\mathrm{Holm}}$ & Sig. \\
\midrule
Q3.1 & 4.62 & 3.93 & -0.519 & $<.001$ & \checkmark \\
Q3.2 & 4.74 & 4.05 & -0.613 & $<.001$ & \checkmark \\
Q3.3 & 4.85 & 4.13 & -0.634 & $<.001$ & \checkmark \\
Q3.4 & 4.80 & 4.05 & -0.660 & $<.001$ & \checkmark \\
Q3.5 & 4.64 & 3.97 & -0.494 & $<.001$ & \checkmark \\
Q3.6 & 4.21 & 3.54 & -0.448 & $<.001$ & \checkmark \\
\bottomrule
\end{tabular}
\end{table}

\begin{table}[ht]
\centering
\caption{Item-level field comparison within the IT-Value family (IT vs.\ IA).}
% ALT: Item-level field comparison within the IT-Value family (IT vs.\ IA), Mann--Whitney $U$ with Holm--Bonferroni family-wise correction ($\alpha=0.05$).
\label{tab:holm_itval}
\small
\begin{tabular}{lrrrrc}
\toprule
Item & IT $M$ & IA $M$ & $r_{rb}$ & $p_{\mathrm{Holm}}$ & Sig. \\
\midrule
Q4.1 & 4.73 & 4.03 & -0.621 & $<.001$ & \checkmark \\
Q4.2 & 4.61 & 3.93 & -0.503 & $<.001$ & \checkmark \\
Q4.3 & 4.69 & 3.94 & -0.553 & $<.001$ & \checkmark \\
Q4.4 & 4.42 & 3.85 & -0.445 & $<.001$ & \checkmark \\
\bottomrule
\end{tabular}
\end{table}

\begin{table}[ht]
\centering
\caption{Item-level field comparison within the Security-Value family (IT vs.\ IA).}
% ALT: Item-level field comparison within the Security-Value family (IT vs.\ IA), Mann--Whitney $U$ with Holm--Bonferroni family-wise correction ($\alpha=0.05$).
\label{tab:holm_secval}
\small
\begin{tabular}{lrrrrc}
\toprule
Item & IT $M$ & IA $M$ & $r_{rb}$ & $p_{\mathrm{Holm}}$ & Sig. \\
\midrule
Q5.1 & 4.47 & 3.76 & -0.564 & $<.001$ & \checkmark \\
Q5.2 & 4.43 & 3.78 & -0.444 & $<.001$ & \checkmark \\
Q5.3 & 4.70 & 3.92 & -0.573 & $<.001$ & \checkmark \\
Q5.4 & 4.69 & 3.95 & -0.525 & $<.001$ & \checkmark \\
Q5.5 & 4.64 & 3.77 & -0.637 & $<.001$ & \checkmark \\
Q5.6 & 4.64 & 3.78 & -0.639 & $<.001$ & \checkmark \\
\bottomrule
\end{tabular}
\end{table}


\subsection{Experience Level}

Kruskal--Wallis tests across the four experience cohorts are significant for all three composites ($H(3)$ between 213.1 and 291.2, all $p<.001$; Table~\ref{tab:exp}), with $\eta^2$ between 0.15 and 0.21 (medium to large). The pattern, however, is \textbf{not monotonic} and therefore cannot be summarized as ``experience increases'' or ``experience decreases'' perceived value. Dunn's post-hoc comparisons with Holm adjustment (Table~\ref{tab:dunn}) locate the effect in the 11--15 year cohort, which scores distinctly \emph{lower} than both the 1--5 and 6--10 year cohorts. The effect is most pronounced on Security Value, where the 11--15 year mean (3.77) sits well below the early-career means (4.32--4.35). The 16+ year cohort posts the highest means on the two value composites, while the 6--10 year cohort leads on Awareness. After correction, the 16+ cohort differs significantly from the 11--15 year trough on both value composites ($z=-5.33$ and $z=-5.12$, both $p<.001$) but not on Awareness ($p=.118$), and none of its comparisons against the 1--5 or 6--10 year cohorts reaches significance. At $n=18$, however, these estimates are unstable, and the analysis does not interpret them further.

On the most defensible reading, mid-career respondents (11--15 years) express more reservation than their junior colleagues, while the most senior group is too small to characterize. This non-monotonic pattern, rather than a linear trend, is what the confidence-interval reporting and non-parametric post-hoc testing specified in Chapter~3 were designed to surface (Figure~\ref{fig:exp}). The multivariate analysis below revisits this pattern, where the mid-career trough proves to be a genuine experience-group difference that survives adjustment for professional field.

Instability in the senior estimates bears on how the pattern should be named. Describing it as a trough implies recovery at the senior end, and that recovery rests entirely on the eighteen respondents in the 16+ cohort. Setting that cohort aside, the three well-populated cohorts trace a decline that begins after the 6--10 year group rather than a dip that resolves. Regression evidence below establishes the 11--15 year effect itself robustly, and the analyses that follow rely on that effect alone; the shape of the curve beyond the 11--15 year cohort does not follow from these data.

\begin{table}[ht]
\centering
\caption{Experience-level comparison of the composite means.}
% ALT: Experience-level comparison (composite means by years of experience). Kruskal--Wallis $H$ with $\eta^2$ effect size $\bigl(\eta^2=(H-g+1)/(N-g)\bigr)$. Cell $n$: 566/489/338/18 respectively.
\label{tab:exp}
\small
\begin{tabular}{lccccrrrr}
\toprule
Composite & 1--5 & 6--10 & 11--15 & 16+ & $H$ & df & $p$ & $\eta^2$ \\
\midrule
Awareness & 4.063 & 4.236 & 3.735 & 4.008 & 213.1 & 3 & $<.001$ & 0.149 \\
IT Value & 4.402 & 4.462 & 3.982 & 4.590 & 229.4 & 3 & $<.001$ & 0.161 \\
Security Value & 4.323 & 4.352 & 3.769 & 4.431 & 291.2 & 3 & $<.001$ & 0.205 \\
\bottomrule
\end{tabular}
\end{table}

\begin{table}[ht]
\centering
\caption{Dunn post-hoc comparisons for the experience-level Kruskal--Wallis tests.}
% ALT: Dunn post-hoc pairwise comparisons for the experience-level Kruskal--Wallis tests, with Holm--Bonferroni adjustment within each composite family. $z$ is the standardized rank-mean difference; a positive $z$ favors the first-named cohort.
\label{tab:dunn}
\small
\begin{tabular}{llrrc}
\toprule
Composite & Cohort comparison & $z$ & $p_{\mathrm{Holm}}$ & Sig. \\
\midrule
Awareness & 1--5 vs 6--10 & -6.09 & $<.001$ & \checkmark \\
 & 1--5 vs 11--15 & 9.51 & $<.001$ & \checkmark \\
 & 1--5 vs 16+ & 0.82 & .411 &  \\
 & 6--10 vs 11--15 & 14.56 & $<.001$ & \checkmark \\
 & 6--10 vs 16+ & 2.39 & .051 &  \\
 & 11--15 vs 16+ & -1.89 & .118 &  \\
\addlinespace
IT Value & 1--5 vs 6--10 & -2.20 & .084 &  \\
 & 1--5 vs 11--15 & 12.52 & $<.001$ & \checkmark \\
 & 1--5 vs 16+ & -1.79 & .146 &  \\
 & 6--10 vs 11--15 & 14.09 & $<.001$ & \checkmark \\
 & 6--10 vs 16+ & -1.22 & .221 &  \\
 & 11--15 vs 16+ & -5.33 & $<.001$ & \checkmark \\
\addlinespace
Security Value & 1--5 vs 6--10 & -0.80 & 1.000 &  \\
 & 1--5 vs 11--15 & 15.07 & $<.001$ & \checkmark \\
 & 1--5 vs 16+ & -0.84 & 1.000 &  \\
 & 6--10 vs 11--15 & 15.35 & $<.001$ & \checkmark \\
 & 6--10 vs 16+ & -0.63 & 1.000 &  \\
 & 11--15 vs 16+ & -5.12 & $<.001$ & \checkmark \\
\bottomrule
\end{tabular}
\end{table}


\begin{figure}[H]
  \centering
    \includegraphics[width=0.82\linewidth, alt={Chart of the three composite means across the four experience cohorts. The pattern is non-monotonic: the 11 to 15 year cohort dips below the others, most clearly on Security Value, where its mean of 3.77 sits below the early-career means of 4.32 to 4.35. The 16+ cohort (n=18) is shown for completeness.}]{figures/fig5_experience.pdf}
  \caption{Composite means by experience cohort.}
  \label{fig:exp}
\end{figure}

\subsection{Association Analysis}

Chi-square tests of independence (Table~\ref{tab:assoc}) are consistent with the associations anticipated in Chapter~3. Professional field is associated with prior exposure (Q6.3$\times$Q1.2, Cram\'er's $V=0.21$), echoing the familiarity gap above: IT respondents are more likely to report recent exposure than IA/Cyber respondents. The field-by-exposure test excludes the Other field group and the two respondents without an exposure response, yielding a three-by-two table with $df=2$ and $n=1{,}339$; Appendix~B reports the full contingency table. Experience is only weakly associated with IT investment willingness ($V=0.05$) and with security investment willingness ($V=0.03$). The latter is unsurprising given that Q5.7 is almost invariant and so has little association to detect with any variable.

\begin{table}[ht]
\centering
\caption{Chi-square tests of independence with Cram\'er's $V$.}
% ALT: Chi-square tests of independence with Cram\'er's $V$. The field $\times$ exposure test excludes the Other field group and two respondents without an exposure response (three-by-two table, hence $df=2$ and $n=1{,}339$).
\label{tab:assoc}
\small
\begin{tabular}{lrrrrr}
\toprule
Association & $\chi^2$ & df & $p$ & Cram\'er's $V$ & $n$ \\
\midrule
Field $\times$ prior exposure (Q6.3$\times$Q1.2) & 60.3 & 2 & $<.001$ & 0.212 & 1{,}339 \\
Experience $\times$ IT investment (Q6.4$\times$Q4.5) & 6.0 & 6 & .421 & 0.046 & 1{,}411 \\
Experience $\times$ Sec investment (Q6.4$\times$Q5.7) & 2.2 & 6 & .902 & 0.028 & 1{,}411 \\
\bottomrule
\end{tabular}
\end{table}


\section{Focused Breakout Analysis: Professional Field and Experience}

Following the proposal, the comparative analysis above contrasted Information Technology against Information Assurance/Cyber Security\footnote{The field category is labeled ``Information Assurance/Cyber Security'' throughout this chapter following the instrument's combined response option (Chapter~3), which merged into one category the governance sense that Chapter~1 designates Information Assurance and the operational sense it designates cybersecurity. The category therefore spans both, and the anonymous instrument cannot separate a governance-oriented from an operationally oriented respondent within it; Chapter~5 carries this heterogeneity into the interpretation of the field gap.} on the composite scores. This section deepens that analysis along the two demographic axes the instrument captured, in two registers. The descriptive breakouts resolve all four professional-field groups (Information Technology, Information Assurance/Cyber Security, Engineering, and Other) and all four experience cohorts, so that the structure within each domain is visible rather than collapsed into a binary. The inferential \emph{stratified} models, by contrast, cover only the two focal groups, Information Technology and Information Assurance/Cyber Security, consistent with the proposal's design. Engineering and Other enter the confirmatory multivariate regression below, with Information Assurance as the reference category, but are not the subject of the stratified within-domain breakouts. This division tests whether the field and experience effects are confounded and finds them largely independent.

\subsection{Composite Scores Across All Four Field Groups}

Resolving the composites across all four field groups (Table~\ref{tab:field_breakout}, Figure~\ref{fig:field_comp}) shows that the Information Technology--Information Assurance contrast above is the widest separation in the data but not the whole picture. Information Technology posts the highest mean on every composite (4.38 awareness, 4.67 IT Value, 4.57 Security Value), while Information Assurance/Cyber Security is the lowest on every composite (3.73, 3.98, 3.82). The Engineering group sits between the two on all three constructs (4.04, 4.34, 4.24). The Other group falls just below Engineering on every composite (3.92, 4.23, 4.14), placing it third of the four groups and above only Information Assurance/Cyber Security. The Other group is small ($n=70$), heterogeneous, and almost entirely early-career (Table~\ref{tab:field_exp_xtab}). Because it does not stand out on any composite, the analysis does not over-interpret it here; Chapter~5 notes it only as a candidate for focused qualitative study. The whole-sample means in the final column show Information Assurance/Cyber Security below the overall average on every composite while Information Technology sits above it.

Familiarity resolves differently from the composites and merits separate comment, because the Engineering group functions as a natural reference point for the field comparison. On Q1.1, Engineering respondents average 1.70 against 2.88 for Information Technology and 1.20 for Information Assurance/Cyber Security, and 51.9\% of Engineering respondents selected ``Not familiar.'' On the composites Engineering sits between the two focal groups, but on named-discipline familiarity it resembles Information Assurance far more closely than Information Technology. That ordering runs against what the discipline's provenance would predict, since Digital Engineering originates in systems engineering practice and its constituent methods are systems engineering methods. Two readings fit the result. Respondents choosing ``Engineering'' were not restricted to systems engineering, so the option would have admitted software, network, industrial, and other specialties for which the discipline is genuinely unfamiliar. This pattern also matches the umbrella-term argument advanced throughout this dissertation, under which practitioners work with model-based methods, digital twins, and lifecycle management individually while not recognizing the collective label the defense community uses to organize them. On that second reading the Engineering group corroborates the central finding rather than complicating it, because it displays the same separation the sample as a whole displays: capability endorsement of 4.43 on the Section~2 items against named-discipline familiarity of 1.70. Because the instrument recorded no engineering subdiscipline, it cannot distinguish the two readings, and Chapter~5 carries the ambiguity into the interpretation of the field comparison.

\begin{table}[ht]
\centering
\caption{Composite scores by professional field (all four groups).}
% ALT: Composite scores by professional field (all four groups). Values are composite means on the 1--5 scale with 95\% confidence intervals; the final column is the whole-sample mean for reference.
\label{tab:field_breakout}
\small
\setlength{\tabcolsep}{4pt}
\begin{tabular}{@{}lccccccccc@{}}
\toprule
 & \multicolumn{2}{c}{IT} & \multicolumn{2}{c}{IA/Cyber} & \multicolumn{2}{c}{Engineering} & \multicolumn{2}{c}{Other} & Overall \\
\cmidrule(lr){2-3}\cmidrule(lr){4-5}\cmidrule(lr){6-7}\cmidrule(lr){8-9}
Composite & $M$ & 95\% CI & $M$ & 95\% CI & $M$ & 95\% CI & $M$ & 95\% CI & $M$ \\
\midrule
Awareness & 4.38 & [4.34, 4.41] & 3.73 & [3.70, 3.77] & 4.04 & [4.00, 4.08] & 3.92 & [3.82, 4.01] & 4.04 \\
IT Value & 4.67 & [4.64, 4.70] & 3.98 & [3.94, 4.01] & 4.34 & [4.30, 4.38] & 4.23 & [4.14, 4.32] & 4.33 \\
Security Value & 4.57 & [4.54, 4.60] & 3.82 & [3.78, 3.86] & 4.24 & [4.20, 4.29] & 4.14 & [4.04, 4.25] & 4.20 \\
\bottomrule
\end{tabular}

\smallskip
\footnotesize Group sizes: IT $n=501$, IA/Cyber $n=501$, Engineering $n=339$, Other $n=70$.
\end{table}


\begin{figure}[H]
  \centering
    \includegraphics[width=0.95\linewidth, alt={Grouped bar chart of the three composite means for all four professional-field groups with 95 percent confidence intervals. Information Assurance/Cyber Security is the lowest group on every composite (3.73, 3.98, 3.82), while Information Technology is highest on every composite, including awareness (4.38), IT Value (4.67), and Security Value (4.57).}]{figures/fig7_field_composites.pdf}
  \caption{Composite means by professional field across all four groups, with 95\% confidence intervals.}
  \label{fig:field_comp}
\end{figure}

The same ordering appears in item-block agreement (Table~\ref{tab:field_section_agree}). Information Assurance/Cyber Security records markedly lower agreement on every block, most starkly on the Security Value items (66.8\% versus 94.7\% for Information Technology and 82.6\% for Other), the domain in which one might most expect security professionals to perceive value. That the security audience expresses the lowest agreement on security-value items is the most counterintuitive descriptive result in the breakout. Because this group mixes governance-oriented and operationally oriented respondents that the anonymous instrument cannot separate, two readings are available, and the data cannot decide between them. The first is an awareness deficit specific to this community. The second is a valid relevance judgment by operationally oriented respondents that a model- and document-centric paradigm fits hands-on defensive work less directly than it fits systems or lifecycle engineering. Chapter~5 takes up both readings and their differing implications for targeted professional development.

\begin{table}[ht]
\centering
\caption{Item-block agreement (percentage of responses at score $\geq 4$) by professional field.}
% ALT: Item-block agreement by professional field. Each cell is the percentage of responses at score $\geq 4$ (agree or strongly agree) for that Section 2--5 item block, by field, with a whole-sample overall column.
\label{tab:field_section_agree}
\small
\begin{tabular}{lccccc}
\toprule
Item block & IT & IA/Cyber & Engineering & Other & Overall \\
\midrule
Awareness (S2) & 98.3 & 88.1 & 96.3 & 92.6 & 93.9 \\
Applicability (S3) & 94.7 & 76.2 & 88.9 & 86.6 & 86.3 \\
IT Value (S4) & 96.3 & 76.7 & 91.1 & 88.6 & 87.7 \\
Security Value (S5) & 94.7 & 66.8 & 86.6 & 82.6 & 82.2 \\
\bottomrule
\end{tabular}
\end{table}


\subsection{Experience Composition Within Each Field}

Experience composition is largely independent of professional field; the two are only weakly associated (Table~\ref{tab:field_exp_xtab}, Figure~\ref{fig:field_exp}). The cross-tabulation is computed on the full sample. The association is small ($\chi^2(9)=9.9$, Cram\'er's $V=0.05$), and the four field groups share broadly the same experience profile. Each concentrates in the two earliest cohorts: Information Technology places 40.3\% at 1--5 years and 35.5\% at 6--10 years, Information Assurance/Cyber Security 37.7\% and 36.3\%, Engineering 42.5\% at 1--5 years and 33.0\% at 6--10 years, and the Other group is the most early-career (44.3\% at 1--5 years). This shared experience profile removes the confound the headline comparisons might otherwise carry.

\begin{table}[ht]
\centering
\caption{Professional field $\times$ experience cross-tabulation.}
% ALT: Professional field $\times$ experience cross-tabulation. Cells show respondent counts with within-field row percentages; the final column is the field total.
\label{tab:field_exp_xtab}
\small
\begin{tabular}{lccccc}
\toprule
Field & 1--5 yrs & 6--10 yrs & 11--15 yrs & 16+ yrs & Total \\
\midrule
IT & 202 (40.3\%) & 178 (35.5\%) & 112 (22.4\%) & 9 (1.8\%) & 501 \\
IA/Cyber & 189 (37.7\%) & 182 (36.3\%) & 123 (24.6\%) & 7 (1.4\%) & 501 \\
Engineering & 144 (42.5\%) & 112 (33.0\%) & 81 (23.9\%) & 2 (0.6\%) & 339 \\
Other & 31 (44.3\%) & 17 (24.3\%) & 22 (31.4\%) & 0 (0.0\%) & 70 \\
\bottomrule
\end{tabular}

\smallskip
\footnotesize $\chi^2(9)=9.9$, $p=.355$, Cram\'er's $V=0.05$, $n=1{,}411$.
\end{table}


\begin{figure}[H]
  \centering
    \includegraphics[width=0.92\linewidth, alt={Stacked bar chart of the experience-cohort composition within each professional field, as within-field percentages. The four groups occupy broadly similar profiles: each concentrates in the two earliest cohorts, with Information Technology at 40.3 percent in 1 to 5 years and 35.5 percent in 6 to 10 years, Information Assurance/Cyber Security at 37.7 and 36.3 percent, Engineering at 42.5 and 33.0 percent, and the Other group the most early-career at 44.3 percent in 1 to 5 years.}]{figures/fig8_field_experience.pdf}
  \caption{Experience composition within each professional field (within-field percentages).}
  \label{fig:field_exp}
\end{figure}

\subsection{Experience Deviation from the Overall Averages}

Examining how each experience cohort deviated from the whole-sample average sharpens the non-monotonic pattern reported in the comparative analysis (Table~\ref{tab:exp_deviation}, Figure~\ref{fig:exp_dev}). Relative to the grand mean, the 11--15 year cohort sits below average on all three composites by a consistent margin, at $-0.31$ on Awareness, $-0.34$ on IT Value, and $-0.43$ on Security Value. The 6--10 year cohort sits above average on all three, at $+0.19$, $+0.14$, and $+0.15$ respectively. The 1--5 year cohort stays close to the overall average throughout. The 16+ year cohort posts positive deviations on the two value composites but, at $n=18$, is statistically unstable and appears for completeness only. The deviation view makes explicit that the experience effect is a mid-career trough at 11--15 years rather than a monotonic increase or decrease across the career span.

\begin{table}[ht]
\centering
\caption{Experience-level deviation from the overall mean.}
% ALT: Experience-level deviation from the overall mean. Each cell gives the cohort composite mean with its deviation from the whole-sample grand mean in parentheses.
\label{tab:exp_deviation}
\small
\begin{tabular}{lccccc}
\toprule
Composite & Overall $M$ & 1--5 yrs & 6--10 yrs & 11--15 yrs & 16+ yrs \\
\midrule
Awareness & 4.04 & 4.06 (+0.02) & 4.24 (+0.19) & 3.73 (-0.31) & 4.01 (-0.04) \\
IT Value & 4.33 & 4.40 (+0.08) & 4.46 (+0.14) & 3.98 (-0.34) & 4.59 (+0.27) \\
Security Value & 4.20 & 4.32 (+0.12) & 4.35 (+0.15) & 3.77 (-0.43) & 4.43 (+0.23) \\
\bottomrule
\end{tabular}

\smallskip
\footnotesize Cohort $n$: 566 / 489 / 338 / 18. The 16+ cohort ($n=18$) is shown for completeness but is statistically unstable.
\end{table}


\begin{figure}[H]
  \centering
    \includegraphics[width=0.95\linewidth, alt={Bar chart of each experience cohort's deviation from the grand composite mean. The 11 to 15 year cohort is below average on all three composites (Awareness -0.31, IT Value -0.34, Security Value -0.43), the 6 to 10 year cohort is above average on all three (+0.19, +0.14, +0.15), the 1 to 5 year cohort sits close to the overall average, and the 16+ cohort (n=18) posts positive deviations but is statistically unstable.}]{figures/fig9_exp_deviation.pdf}
  \caption{Deviation of each experience cohort from the overall composite mean.}
  \label{fig:exp_dev}
\end{figure}

\subsection{Focused Models: Disentangling Field from Experience}

Together, the two preceding subsections show that the mid-career trough and the Information Technology--Information Assurance gap are largely independent findings rather than a single confounded one. Because professional field and experience are only weakly associated, the 11--15 year cohort is not disproportionately Information Assurance, so neither effect is reducible to the other. Resolving the composites by experience \emph{within} the two focal domains corroborates this separation (Table~\ref{tab:exp_field_focus}).

Two results follow. First, the 11--15 year cohort is a genuine local low point within \emph{both} focal domains: relative to its own domain mean, it deviates by $-0.30$ to $-0.38$ within Information Technology and by $-0.31$ to $-0.44$ within Information Assurance/Cyber Security across the three composites. The mid-career reservation is therefore a real experience-group pattern present in each community rather than an artifact of field composition. Second, the Information Technology--Information Assurance gap is \emph{not} an artifact of experience: within each of the three interpretable experience cohorts (1--5, 6--10, and 11--15 years), Information Technology exceeds Information Assurance on every composite. On awareness, for example, the contrast runs 4.41 versus 3.72 at 1--5 years, 4.53 versus 3.97 at 6--10 years, and 4.08 versus 3.42 at 11--15 years. Information Technology remains higher on all three composites even in the 16+ year cohort, though with only 9 Information Technology respondents against 7 Information Assurance respondents, that cohort is too small to interpret. Holding experience constant leaves field the dominant axis of variation, and the mid-career trough an independent experience-group difference within each field. These focused models, restricted as the design specifies to the Information Technology and Information Assurance/Cyber Security groups, qualify the comparative results without overturning them.

\begin{table}[ht]
\centering
\caption{Experience-resolved composite means within the IT and IA/Cyber domains.}
% ALT: Experience-resolved composite means within the two focal domains (IT and IA/Cyber). Each cell gives the within-domain cohort mean with its deviation from that domain's own mean in parentheses.
\label{tab:exp_field_focus}
\small
\setlength{\tabcolsep}{5pt}
\begin{tabular}{@{}llccccc@{}}
\toprule
Composite & Domain & Domain $M$ & 1--5 yrs & 6--10 yrs & 11--15 yrs & 16+ yrs \\
\midrule
Awareness & IT & 4.38 & 4.41 (+0.03) & 4.53 (+0.15) & 4.08 (-0.30) & 4.41 (+0.03) \\
Awareness & IA & 3.73 & 3.72 (-0.01) & 3.97 (+0.24) & 3.42 (-0.31) & 3.53 (-0.20) \\
\addlinespace
IT Value & IT & 4.67 & 4.74 (+0.07) & 4.77 (+0.10) & 4.36 (-0.31) & 4.85 (+0.18) \\
IT Value & IA & 3.98 & 4.03 (+0.05) & 4.16 (+0.18) & 3.63 (-0.35) & 4.21 (+0.23) \\
\addlinespace
Security Value & IT & 4.57 & 4.66 (+0.09) & 4.68 (+0.11) & 4.19 (-0.38) & 4.69 (+0.12) \\
Security Value & IA & 3.82 & 3.93 (+0.11) & 4.00 (+0.18) & 3.38 (-0.44) & 4.02 (+0.20) \\
\bottomrule
\end{tabular}
\end{table}


\subsection{Multivariate Corroboration}

Stratified breakouts suggest, but do not formally test, that professional field and experience each contribute independently to the value differences. An ordinary least squares regression of each value composite on field, experience, and self-reported familiarity makes that test explicit~\parencite{Cohen_Cohen_West_Aiken_2003}. The models are fit on all 1{,}411 respondents. They account for a substantial share of variance (IT~Value $R^2=0.64$; Security~Value $R^2=0.67$; Table~\ref{tab:regression}) and yield three results that sharpen the descriptive picture. Because the value composites are bounded on the 1--5 scale and left-skewed, a robustness check refit the models with HC3 heteroscedasticity-robust standard errors. The refit left every coefficient's direction and significance intact (Appendix~B), so the conclusions do not depend on the homoscedastic ordinary-least-squares assumption. Because professional field and self-reported familiarity are themselves correlated, the analysis checked collinearity directly: the maximum variance inflation factor is approximately 2.0 (familiarity 1.58, Information Technology 2.00), well below any conventional threshold, so the field and familiarity coefficients are separately identified.

Professional field remains a strong predictor once the models hold experience and familiarity constant. Relative to Information Assurance/Cyber Security, Information Technology respondents score $0.46$ points higher on IT~Value and $0.50$ higher on Security~Value, both at $p<.001$. The Engineering and Other groups also score higher than Information Assurance ($0.30$ and $0.35$ for Engineering, $0.23$ and $0.31$ for Other), likewise at $p<.001$. The field gap therefore survives adjustment for experience and familiarity, supporting the within-cohort pattern above.

Most consequentially, the mid-career experience effect persists once the models adjust for field and familiarity: relative to the 1--5 year reference, the 11--15 year cohort scores $0.402$ points lower on IT~Value and $0.534$ lower on Security~Value, both at $p<.001$. This mirrors the lowest raw means the cohort posted in the omnibus analysis. This is consistent with a genuine experience-group difference rather than a composition artifact: with the two variables as nearly orthogonal as the cross-tabulation showed, the mid-career reservation cannot be explained away by field membership. This is the inferential counterpart to the within-domain trough identified earlier.

Familiarity, finally, is a positive predictor of perceived value (IT~Value $b=0.132$, standardized $\beta=0.34$; Security~Value $b=0.142$, $\beta=0.33$; both $p<.001$; standardized coefficients from the same models, not shown in Table~\ref{tab:regression}).
% Verified: standardized beta = 0.34 (IT Value) and 0.33 (Security Value) reproduced in Analysis/DigitalEngineering_Analysis.ipynb (Section 11e, standardized-coefficients cell).
The result qualifies, without contradicting, the headline finding that respondents perceive value despite low familiarity: across respondents greater familiarity accompanies greater perceived value, yet the overall value means remain high because even unfamiliar respondents endorse the capability descriptions. Rank-order correlations point the same way: familiarity correlates moderately with both value composites (Spearman rank-order correlations $\rho=0.63$ with IT~Value and $\rho=0.62$ with Security~Value, both $p<.001$)~\parencite{Spearman_1904}.
% Verified: Spearman rho = 0.632, 0.616, 0.734, 0.731 (all p < .001, n = 1,411) reproduced in Analysis/DigitalEngineering_Analysis.ipynb (Section 11f).

\begin{table}[ht]
\centering
\caption{Regression of the value composites on field, experience, and familiarity.}
% ALT: Ordinary least squares regression of the value composites on professional field, experience level, and self-reported familiarity. Coefficients are unstandardized, with standard errors in parentheses; reference categories are Information Assurance/Cyber Security (field) and 1--5 years (experience).
\label{tab:regression}
\small
\begin{tabular}{lcc}
\toprule
Predictor & IT Value & Security Value \\
\midrule
Familiarity (Q1.1) & $0.132^{***}$ (0.008) & $0.142^{***}$ (0.008) \\
Field: Information Technology & $0.458^{***}$ (0.022) & $0.496^{***}$ (0.023) \\
Field: Engineering & $0.301^{***}$ (0.020) & $0.350^{***}$ (0.021) \\
Field: Other & $0.232^{***}$ (0.036) & $0.305^{***}$ (0.038) \\
Experience: 6--10 years & $0.082^{***}$ (0.017) & $0.055^{**}$ (0.018) \\
Experience: 11--15 years & $-0.402^{***}$ (0.019) & $-0.534^{***}$ (0.020) \\
Experience: 16+ years & $0.168^{*}$ (0.067) & $0.092$ (0.070) \\
Intercept & $3.888^{***}$ (0.018) & $3.759^{***}$ (0.019) \\
\midrule
$R^2$ & 0.635 & 0.672 \\
Adjusted $R^2$ & 0.633 & 0.670 \\
$N$ & 1{,}411 & 1{,}411 \\
\bottomrule
\end{tabular}

\smallskip
\footnotesize $^{*}p<.05$, \;$^{**}p<.01$, \;$^{***}p<.001$. Reference categories: Information Assurance/Cyber Security; 1--5 years of experience.
\end{table}


\section{Bias Diagnostics}

The methodology specified two planned diagnostics for response bias: a wave (non-response) analysis and a priming diagnostic. Both diagnostics ran as planned.

\subsection{Wave (Non-Response) Analysis}

Splitting the sample at the median submission time into early- and late-arriving halves reveals no meaningful differences on any of the three composites. Every effect size is negligible (Cohen's $d$ between $-0.02$ and $-0.04$) and none approaches significance (all Mann--Whitney $p>.5$), with late-arriving respondents scoring only marginally higher than early-arriving ones (Table~\ref{tab:wave}). The split is by date of submission rather than completion duration, following the early-versus-late tradition of non-response sensitivity analysis~\parencite{Armstrong_Overton_1977}. Under that logic, in which late respondents proxy for non-respondents, the near-identical composite means across the two halves indicate that late arrivals differ negligibly from early arrivals, so the value composites show little sign of non-response bias. The null wave effect is reassuring on its own terms. It does not, however, by itself dispose of the recruitment-frame concern developed in Chapter~5, because both halves arrived through the same self-selected channels rather than from a random draw of the profession; Chapter~5 carries that concern into the limitations alongside this null result.

\begin{table}[ht]
\centering
\caption{Wave (non-response) analysis: early- vs.\ late-arriving respondents.}
% ALT: Wave (non-response) analysis: early- vs.\ late-arriving half by submission time (Armstrong--Overton early/late logic). $p$-values from the Mann--Whitney $U$ test; $d$ = Cohen's $d$, with a negative $d$ indicating the late-arriving half scores higher.
\label{tab:wave}
\small
\begin{tabular}{lrrrr}
\toprule
Composite & Early $M$ & Late $M$ & $p$ & $d$ \\
\midrule
Awareness & 4.033 & 4.054 & .523 & -0.04 \\
IT Value & 4.321 & 4.328 & .802 & -0.02 \\
Sec Value & 4.192 & 4.212 & .527 & -0.04 \\
\bottomrule
\end{tabular}
\end{table}


\subsection{Priming Diagnostic}

The priming check specified in Chapter~3 compares the value composites between high-familiarity respondents (Q1.1 $=4,5$) and low-familiarity respondents (Q1.1 $=1,2$). High-familiarity respondents report substantially higher perceived value on both composites, with differences of 0.54 scale points on IT Value and 0.57 on Security Value and $d\approx1.3$ (Table~\ref{tab:priming}). By the logic stated in Chapter~3, this weighs \emph{against} the pure-priming hypothesis. If the contextual descriptions in the question stems were manufacturing agreement independent of knowledge, the two familiarity groups would converge. Here they diverge sharply, a pattern more consistent with genuine perception differences informed by prior knowledge (Figure~\ref{fig:prime}). The corollary is that aggregate value estimates are sensitive to the familiarity composition of the sample: because this sample is predominantly low-familiarity, the overall value means rest largely on the low-familiarity stratum, with the high-familiarity stratum reporting markedly higher value. Chapter~5 weighs what this implies for the broader population.

\begin{table}[ht]
\centering
\caption{Priming diagnostic: value composites by familiarity level.}
% ALT: Priming diagnostic: value composites for high- vs.\ low-familiarity respondents. $p$-values from the Mann--Whitney $U$ test; $d$ = Cohen's $d$.
\label{tab:priming}
\small
\begin{tabular}{lrrrrr}
\toprule
Composite & High $M$ & Low $M$ & $\Delta$ & $p$ & $d$ \\
\midrule
IT Value & 4.751 & 4.212 & 0.539 & $<.001$ & 1.30 \\
Sec Value & 4.654 & 4.081 & 0.573 & $<.001$ & 1.25 \\
\bottomrule
\end{tabular}
\end{table}


\begin{figure}[H]
  \centering
    \includegraphics[width=0.66\linewidth, alt={Bar chart comparing the two value composites between high-familiarity respondents (Q1.1 = 4 or 5) and low-familiarity respondents (Q1.1 = 1 or 2). High-familiarity respondents report higher perceived value on both composites, by 0.54 scale points on IT Value and 0.57 on Security Value, with Cohen's d of about 1.3.}]{figures/fig6_priming.pdf}
  \caption{Value composites by familiarity stratum.}
  \label{fig:prime}
\end{figure}

\section{Supplementary Post-Hoc Robustness and Validity Analyses}

Analyses in this section were not part of the approved proposal's analysis plan (Chapter~3, Research Propositions and Planned Analyses). The candidate added them after data collection concluded to strengthen the study's inferential foundations, a self-initiated step informed by a review of methodological practice in comparable dissertations. They test the robustness and measurement properties of the findings already reported rather than evaluate new propositions. Because they are post-hoc in this sense, the section reports them as corroboration and diagnostic evidence rather than pre-specified confirmatory tests, and each table is annotated accordingly. Their shared purpose is to examine whether the study's central results, and the professional-field contrast in particular, withstand scrutiny the proposal did not anticipate.

\subsection{Common-Method Variance}

Because every measure derives from a single self-report instrument completed at one sitting, common-method variance could in principle inflate the observed associations~\parencite{Podsakoff_2003}. Harman's single-factor test offers a first screen: it asks whether one unrotated factor accounts for most of the variance across all items. It does not. A single factor explains 44.0\% of the variance across the 23 composite items, below the 50\% level conventionally read as evidence of a dominant method factor, and two factors exceed the eigenvalue-one criterion (Table~\ref{tab:cmv}). Harman's test is a lower-bound diagnostic rather than a decisive control, and 44.0\% sits modestly below the threshold, so the conclusion is narrow: a single method factor does not dominate these data.

Beyond Harman's screen, the analysis considered a fuller latent-method-factor control and reports it here as a null attempt rather than a result. Fitting an unmeasured latent method construct, in which a common factor loads on all items alongside the substantive factors, is uninformative here because the two value constructs are nearly collinear (heterotrait--monotrait ratio $0.99$, reported below). Under such low discriminant validity the added factor absorbs genuine shared substantive variance rather than isolating method variance, a documented failure mode of the technique~\parencite{Podsakoff_2003}. The apparent method variance it returns is therefore not a valid estimate, and the chapter does not report it as one.

The decisive point is conceptual as well as statistical. A method bias shared identically by every respondent shifts both groups equally, so it cannot by itself manufacture the between-groups mean difference that carries the study's central inference, the large Information Technology versus Information Assurance gap. A uniform response style is therefore an unlikely explanation for that gap. A response style that differs \emph{by group}, however, such as more conservative endorsement among security professionals, is a group-correlated method effect Harman's test does not exclude. The differential-item-functioning screen below addresses exactly this, showing that the items function equivalently across the two communities once respondents are matched on their own composite score. That match is itself imperfect, because the composite carries any uniform style inflation. Three strands of evidence therefore converge. Harman's test sets a lower bound, the between-groups logic rules out a uniform bias, and the measurement-equivalence screen shows the items behaving alike across the two communities. Taken together they make common-method variance an unlikely sole explanation for the field and familiarity relationships, although a differential response style remains among the accounts Chapter~5 weighs.

\begin{table}[htbp]\centering
\caption{Common-method variance: Harman's single-factor test.}\label{tab:cmv}
% ALT: Common-method variance check (Harman's single-factor test): items entered, factors with eigenvalue above one, and variance explained by the first unrotated factor against the 50 percent dominance threshold.
\begin{tabular}{@{}p{3.6in}r@{}}\toprule
\textbf{Diagnostic} & \textbf{Value} \\ \midrule
Items entered (all composite items) & 23 \\
Factors with eigenvalue $>1$ & 2 \\
Variance explained by first unrotated factor & 44.0\% \\
Single-factor dominance threshold & 50\% \\ \bottomrule
\end{tabular}

\smallskip
\parbox{3.9in}{\footnotesize \textit{Note.} Post-hoc analysis, not specified in the proposal. A first factor below 50\% indicates common-method variance is unlikely to account for the observed relationships.}
\end{table}

\subsection{Construct Validity of the Composite Scores}

The comparative analyses treat the three composites as summative indices of related perceptions. A confirmatory three-factor model checks their measurement properties more strictly, estimating composite reliability and average variance extracted for each~\parencite{Fornell_Larcker_1981}. Composite reliability is high throughout, at 0.86 for Awareness, 0.87 for IT Value, and 0.87 for Security Value (Table~\ref{tab:construct}), reinforcing the Cronbach's $\alpha$ values in Table~\ref{tab:comp}. Convergent validity, indexed by average variance extracted, falls just below the 0.50 benchmark for all three composites (0.47 for Awareness, 0.46 for IT Value, and 0.45 for Security Value), indicating that items within each composite share somewhat less than half their variance with the intended construct. Discriminant validity is not established under the Fornell--Larcker criterion: the three composites intercorrelate between 0.74 and 0.85, exceeding the square roots of their average variance extracted (0.67 to 0.69).

The heterotrait--monotrait ratio of correlations sharpens this picture and is more informative than the Fornell--Larcker test here~\parencite{Henseler_Ringle_Sarstedt_2015}. Awareness is discriminant from both value composites, with HTMT ratios of 0.83 against each, below the conventional 0.85 threshold; the two value composites, by contrast, are not discriminant from each other (HTMT~$=0.99$), so IT Value and Security Value are empirically a single perceived-value construct rather than two. The confirmatory three-factor model itself fits the data well (CFI~$=0.966$, RMSEA~$=0.040$), so the measurement structure is sound. The discriminant criteria flag the near-identity of the two value dimensions, a separability issue that leaves model fit intact; the analysis accordingly treats the two as related summative indices and never as independent latent variables.

\begin{table}[htbp]\centering
\caption{Construct validity: composite reliability, average variance extracted, and HTMT discriminant ratios.}\label{tab:construct}
% ALT: Construct validity table: composite reliability (CR), average variance extracted (AVE), and root-AVE per composite, plus heterotrait-monotrait (HTMT) ratios per composite pair against the 0.85 discriminant threshold.
\begin{tabular}{@{}lrrr@{}}\toprule
\textbf{Composite} & \textbf{CR} & \textbf{AVE} & \textbf{$\sqrt{\text{AVE}}$} \\ \midrule
Awareness & 0.858 & 0.471 & 0.686 \\
IT Value & 0.868 & 0.455 & 0.675 \\
Sec Value & 0.867 & 0.453 & 0.673 \\
\bottomrule
\end{tabular}
\smallskip

\begin{tabular}{@{}lrc@{}}\toprule
\textbf{Composite pair} & \textbf{HTMT} & \textbf{$<0.85$?} \\ \midrule
Awareness ~ IT Value & 0.830 & yes \\
Awareness ~ Sec Value & 0.830 & yes \\
IT Value ~ Sec Value & 0.987 & no \\
\bottomrule
\end{tabular}

\smallskip
\parbox{4.8in}{\footnotesize \textit{Note.} Post-hoc analysis, not specified in the proposal. CR $\geq .70$ and
AVE $\geq .50$ are the conventional thresholds for adequate convergent validity; the heterotrait-monotrait
ratio (HTMT) discriminant threshold is $0.85$. Values are computed from standardized loadings of a
confirmatory three-factor model and from the item correlation matrix.}
\end{table}

These results carry a direct interpretive consequence, and the HTMT ratios locate it precisely. Awareness is empirically separable from perceived value, but the two value dimensions, IT Value and Security Value, are not separable from each other: respondents who value Digital Engineering for IT operations value it for security operations to almost exactly the same degree. This is theoretically unsurprising, since both dimensions tap a single underlying appraisal of the same described capabilities. The value composites are therefore best understood, and used throughout this chapter, as reliable summative indices of one perceived-value construct measured in two domains rather than as two discriminant latent variables. A consequence follows for how the field gap is counted: because IT Value and Security Value index a single construct, their near-identical field gaps ($d=1.94$ on each) constitute one finding rather than two independent confirmations. The Awareness gap ($d=1.64$), meanwhile, leans partly on the cross-scale, DIF-affected familiarity item Q1.1 and is the least clean of the three. This does not weaken the group comparisons, which depend on the reliability of the indices and not on their mutual separability, but it does bound any claim that IT Value and Security Value capture independent dimensions.

One further consequence bears directly on the Research Question~2 and Research Question~3 distinction. Chapter~1 separates general perceived value from applicability to a respondent's own work, a difference that matters for adoption strategy, and the present items do not separate the two empirically. Testing that separation would require distinct item sets written to it, with general-relevance and own-context stems kept apart by construction. Chapter~5 carries this qualification into the limitations.

\subsection{Robustness of the Field Gap: Bootstrap Intervals and the Combined Value Construct}

Two robustness checks test whether the headline professional-field gap depends on how the value composites were built or on the parametric assumptions behind Cohen's $d$. The first pools all sixteen IT-value and security-value items into a single Perceived Value composite, warranted because the discriminant analysis above showed the two dimensions to be one construct. The pooled composite is highly reliable (Cronbach's $\alpha=0.92$), and the Information Technology versus Information Assurance gap on it is if anything larger than on the separate composites (Cohen's $d=2.08$), so the field difference is not an artifact of splitting value into two correlated halves. The second replaces the normal-theory confidence interval for $d$ with a percentile interval from 5{,}000 bootstrap resamples, which makes no distributional assumption about the bounded, left-skewed composite scores. The bootstrap intervals are narrow and sit far above the conventional large-effect threshold of $0.8$ for every composite. Awareness returns $d=1.64$ with a 95\% interval of $[1.50, 1.79]$, IT Value $1.94$ with $[1.78, 2.11]$, Security Value $1.94$ with $[1.79, 2.11]$, and the combined Perceived Value composite $2.08$ with $[1.92, 2.26]$. These results reinforce the non-parametric Mann--Whitney tests already reported (Table~\ref{tab:robustness_extra}, Figure~\ref{fig:forest}). The field gap is therefore robust to the composite specification, to the distributional assumptions of the parametric effect size, and to the choice between parametric and rank-based inference.

\begin{table}[htbp]\centering
\caption{Field-gap effect sizes with bootstrap confidence intervals and the combined value composite.}\label{tab:robustness_extra}
% ALT: Field-gap effect sizes (Cohen's d) with bootstrap 95 percent confidence intervals for each value composite and for the combined value composite, comparing IT and IA/Cyber respondents.
\small
\begin{tabular}{@{}lrrrl@{}}\toprule
\textbf{Composite} & \textbf{IT} & \textbf{IA} & \textbf{Cohen's }$d$ & \textbf{95\% bootstrap CI} \\ \midrule
Awareness & 4.38 & 3.73 & 1.64 & [1.50, 1.79] \\
IT Value & 4.67 & 3.98 & 1.94 & [1.78, 2.11] \\
Security Value & 4.57 & 3.82 & 1.94 & [1.79, 2.11] \\
Perceived Value (combined) & 4.62 & 3.90 & 2.08 & [1.92, 2.26] \\
\bottomrule
\end{tabular}
\smallskip

{\footnotesize \textit{Note.} Post-hoc analysis, not specified in the proposal. Group means for Information
Technology (IT) and Information Assurance (IA); Cohen's $d$ with a percentile 95\% confidence interval from
5{,}000 bootstrap resamples. The Perceived Value composite pools all sixteen IT-value and security-value
items, which the heterotrait--monotrait analysis showed to measure a single construct ($\alpha=0.92$).}
\end{table}

\begin{figure}[htbp]
  \centering
    \includegraphics[width=0.9\linewidth, alt={Forest plot of the Information Technology minus Information Assurance effect size (Cohen's d) for the Awareness, IT Value, Security Value, and combined Perceived Value composites, each with a 95 percent bootstrap confidence interval. All four effect sizes lie between 1.6 and 2.1 with intervals far to the right of the large-effect reference line at 0.8.}]{figures/fig10_forest.pdf}
  \caption{Professional-field effect sizes (IT versus IA) with 95\% bootstrap confidence intervals.}
  \label{fig:forest}
\end{figure}

\subsection{Measurement Invariance Across Professional Field}

Professional field carries the largest effect in the data, and one candidate explanation for it is differential item functioning: the possibility that Information Assurance and Information Technology respondents at the same underlying level of a construct nonetheless answer particular items differently. A regression-based screen, adapting Zumbo's logistic DIF framework to a linear implementation on the item scores, tests this directly by regressing each item on the respondent's own composite score, professional field, and their interaction. Any residual field effect then represents a difference the matched trait does not explain~\parencite{Zumbo_1999}. Twenty-one of the 23 items show no meaningful differential functioning. Two meet the flag criterion (Table~\ref{tab:invariance}): the familiarity item Q1.1, whose matched-level difference is large ($\beta_{\text{IA}}=-0.99$) and largely restates the familiarity gap itself, and the capability item Q2.5, whose effect is small ($\beta_{\text{IA}}=+0.23$). Setting the familiarity item aside, the residual differences among the substantive capability items are small relative to the composite-level field difference, which approaches two standard deviations (Figure~\ref{fig:box}). The screen therefore supports a substantive reading of the field gap. A single capability item functions slightly differently across the two communities and merits cautious interpretation, but item-level measurement differences cannot account for a separation of the magnitude reported earlier. This evidence weakens the differential-item-functioning account specifically, while leaving the substantive awareness-and-exposure account for Chapter~5 to weigh.

\begin{table}[htbp]\centering
\caption{Measurement-invariance screen: regression-based DIF, IT vs.\ IA/Cyber.}\label{tab:invariance}
%% ALT: Measurement-invariance screen using regression-based differential item functioning (DIF) between IT and IA/Cyber respondents: per-item beta, uniform-DIF p-value, and flag status.
\begin{tabular}{@{}lrrc@{}}\toprule
\textbf{Item} & \textbf{$\beta_{\text{IA}}$ (matched)} & \textbf{$p$} & \textbf{DIF flag} \\ \midrule
Q1.1 & -0.99 & $<.001$ & yes \\
Q2.5 & +0.23 & $<.001$ & yes \\
Q5.2 & +0.16 & .003 & no \\
Q2.4 & +0.16 & $<.001$ & no \\
Q2.6 & +0.16 & $<.001$ & no \\
Q2.1 & +0.15 & $<.001$ & no \\
Q2.2 & +0.15 & $<.001$ & no \\
Q2.3 & +0.14 & $<.001$ & no \\
\bottomrule
\end{tabular}

\smallskip
\parbox{4.8in}{\footnotesize \textit{Note.} Post-hoc analysis, not specified in the proposal. $\beta_{\text{IA}}$ is the
IA-vs-IT difference in item response (scale points) after matching on the item's composite score
(a linear implementation of Zumbo's DIF framework). 2 of 23 items met the flag criterion ($p<.001$ and
$|\beta_{\text{IA}}| \geq 0.20$). All flagged effects are small relative to the composite-level
field difference, which approaches two standard deviations.}
\end{table}

\subsection{Sensitivity to Straight-Lining}

The straight-lining check reported in Appendix~B identified 22 respondents (1.6\% of the 1{,}390 with the complete agreement battery) who gave an identical answer to all twenty-two agreement-scale items. Removing them tests whether the high endorsement is an artifact of non-differentiating responding. It is not: with the straight-liners excluded, Section~2 agreement is essentially unchanged at 93.9\%, because so few respondents straight-lined that their removal does not move the aggregate. Acquiescent straight-lining is therefore a negligible presence and does not manufacture the high endorsement.

\subsection{Exposure, Familiarity Composition, and the Field Gap}

Differential prior exposure, rather than field membership, could explain the professional-field difference, because Information Technology respondents were far more likely to report recent exposure to Digital Engineering than Information Assurance respondents (32.1\% versus 11.8\%, a 161-to-59 split among the two focal groups). Adding a recent-exposure indicator to the value regressions of Table~\ref{tab:regression} tests this directly. The field coefficient is essentially unchanged (Information Technology relative to Information Assurance $b=0.45$ for IT~Value and $0.49$ for Security~Value, both $p<.001$, against $0.46$ and $0.50$ without the exposure term), and the exposure term itself is small ($b=0.08$ for both composites, $p<.001$). The stark exposure imbalance across fields therefore does not account for the field gap in perceived value, which persists once exposure is held constant. Because the exposed Information Assurance subgroup numbers 59 respondents, this is reported as a regression adjustment rather than a within-stratum comparison.

Bimodal familiarity, described earlier in this chapter, supports a second and more demanding version of the same test. If the sample comprises a large naive stratum and a smaller well-acquainted one, and if the two focal fields draw unequally from them, then stratum composition rather than professional field could carry the difference. Information Technology and Information Assurance do draw unequally: 206 of the 501 Information Technology respondents fall in the high-familiarity stratum (Q1.1 of 4 or 5) against 11 of the 501 Information Assurance respondents. Adding an indicator for that stratum to the value regressions holds the composition constant. Field coefficients fall from 0.689 to 0.602 on IT Value and from 0.746 to 0.655 on Security Value, retaining 87\% of their magnitude in both cases, and remain significant at $p<.001$ under HC3-robust standard errors. Entering the exposure indicator and the stratum indicator together retains 86\%. The two-stratum structure of the familiarity distribution therefore accounts for roughly one seventh of the field gap and leaves the remainder standing. This analysis was not specified in the proposal and is reported, with the other supplementary analyses in this section, as a self-initiated robustness check. It narrows the recruitment-frame concern without resolving it, because a channel difference that delivered systematically different respondents within the same familiarity stratum would not be visible to this adjustment.

\subsection{Response-Authenticity Screening (Post-Hoc Validation)}\label{sec:forensic-validation}

Because the survey ran anonymously through an open link, the data were exposed to the automated, coordinated, and synthetic responding that increasingly threatens online instruments of this kind~\parencite{Pozzar_Bots_2020, Griffin_Bots_2022, Westwood_LLM_2025}. A growing share of this contamination is coordinated human fraud, such as incentive-farming rings sharing an answer key, rather than crude automation alone~\parencite{Jaffe_NotBots_2026}. That is precisely why a screen keyed on the shared response vector is apt: it detects coordinated duplication whether a script or a coordinated human effort produced it. As a post-hoc validation not specified in the proposal, the analysis added a response-authenticity screen to the pre-registered quality screen described in Chapter~3. This screen and its validation fall outside the original scope of the dissertation's research questions, which concern awareness and perceived value of Digital Engineering rather than survey methodology. The chapter nonetheless develops and reports the screen because it carries independent academic value: it contributes a repeatable, openly released, and validated procedure that researchers across fields conducting anonymous online surveys can adopt, beyond its immediate service to this study's data quality. The screen operates only on the submitted response vectors and their submission timing; the anonymous design assigns every respondent the same placeholder identity, so identity fields carry no forensic information and are not used.

The pre-registered quality screen had already removed 188 rapid or abandoned submissions from the 1{,}600 raw submissions. Screening for authenticity then examined the remaining sample for coordinated duplication using a documented rule: it flagged any complete response vector shared by four or more submissions as a candidate scripted cluster. It also flagged near-duplicate satellites within a small edit distance of such a cluster and the later copies of smaller duplicate pairs. A single cluster stood out unambiguously, and its structure identifies it as an unsophisticated automated attack of the kind often termed a script-kiddie attack: one reusing a single hard-coded payload with no attempt to vary the answers, randomize the submission timing, or disguise the payload. The payload was one fixed response vector, a straight-line pattern of the maximum agreement value on every item paired with a single fixed demographic profile, submitted 71 times. Every one of the 71 submissions carried the byte-for-byte identical response vector, each completed in a median of 1.1 minutes, and the submissions arrived in a burst spaced a median of 1.1 minutes apart. Nothing about the payload changed from one submission to the next, precisely the weakness that makes such an attack trivial to detect: a genuine population never reproduces one complete, fully specified answer vector dozens of times, whereas a script replaying hard-coded answers does so by construction.

The layered screen caught the entire burst through two independent and mutually reinforcing mechanisms. The pre-registered timing screen removed 70 of the 71 submissions on completion speed alone, because their sub-three-minute completions fell below the rapid-response threshold. The response-authenticity screen then recognized all 71 as a single coordinated cluster on the strength of their shared complete response vector, a count far above the four-submission threshold. In doing so it caught the one member of the burst whose completion time had drifted just past the timing cutoff and would otherwise have survived. Together the two screens removed every submission in the attack: the timing filter caught the bulk on speed, and the vector screen recognized the attack as a single object and closed the one-response gap the timing filter left. The burst's net effect on the analytic sample was therefore one additional excluded submission beyond those the timing screen had already caught, and the sample stands at $N=1{,}411$.

The screen also flagged a smaller set of candidate duplicate rows beyond that cluster, and these received individual manual scrutiny rather than automatic exclusion, in the documented and disclosed manner best-practice data screening recommends~\parencite{DeSimone_Screening_2015}. On inspection they did not resemble coordinated responding. Their completion times were normal, with a median near nine minutes rather than the sub-minute times of the automated burst, and the members of each flagged pattern had arrived days or weeks apart rather than in a burst. Their response profiles were also common among clearly authentic respondents: an Information Technology professional agreeing strongly across the value items is one of the most frequent genuine profiles in the sample. Two such respondents independently producing the same response vector is therefore an expected coincidence near the response ceiling rather than evidence of duplication. The analysis judged these cases authentic and retained them. To confirm that this judgment does not shape the results, the analysis repeated every test with the full set of automatically flagged rows removed ($N=1{,}380$). The two samples yield materially identical results. Composite reliabilities move by less than $0.01$. The Information Technology versus Information Assurance effect sizes remain very large, with Cohen's $d$ of $1.58$, $1.88$, and $1.89$ on the fully screened sample against $1.64$, $1.94$, and $1.94$ on the primary one. Every association and diagnostic in this chapter is unchanged in direction and significance. Retaining or excluding the manually scrutinized cases changes no conclusion.

Before applying the screen, the analysis validated it against synthetic data with known ground truth, so that its error characteristics are established rather than assumed~\parencite{Fawcett_ROC_2006, Ilagan_ModelAgnostic_2023, Dupuis_Indices_2019}. That validation, detailed in Appendix~B, characterizes both what the screen catches---coordinated and duplicated submissions---and where it is deliberately conservative: applied mechanically to ceiling-heavy data it would over-exclude genuine strong endorsers at roughly one to two percent, the same borderline rate seen in this sample. For exactly that reason the procedure pairs the automated rule with human adjudication rather than applying it alone. Reported as a data-quality validation, the exercise establishes that a single automated submission burst reached the instrument and was removed, that no further coordinated structure survives in the analytic sample, and that the findings are stable to the treatment of every borderline case the screen identified.

Taken together, the three diagnostics in this section speak to a single question: whether the analyzed responses are organic, the product of real and engaged professional participants rather than of automation, coordination, or inattention. Each addresses a distinct threat, and the evidence converges. The wave analysis found no material difference between early and late respondents on any composite ($|d| \leq 0.04$), so the sample shows no self-selection signal along the order of submission. Straight-lining, the signature of inattentive or acquiescent responding, is present in only $1.6\%$ of the sample, and its removal moves the headline endorsement rates by fractions of a percentage point. The response-authenticity screen isolated the single automated submission burst that reached the instrument, removed it, and confirmed that no further coordinated structure survives among the retained $1{,}411$ responses. None of these results proves that every retained submission is genuine, a claim no screen can support. Two residual limits remain explicit: professional field is self-reported and unverifiable under the anonymous design, and individual misrepresentation of role lies beyond the reach of any content-based screen. Within those bounds, the converging evidence supports reading the analytic sample as a body of organic professional responses, and it establishes the empirical basis on which the findings that follow rest.

\section{Summary of Findings}

Table~\ref{tab:answers}, presented with the direct answers earlier in this chapter, states the finding for each research question. Table~\ref{tab:glance} sets the four propositions against the evidence before the narrative summary that follows.

\begin{table}[htbp]
  \centering
  \caption{The four propositions against the evidence.}\label{tab:glance}
  % ALT: Table summarizing the four research propositions against the findings and the key evidence supporting each.
  \small
  \begin{tabular}{@{}>{\raggedright\arraybackslash}p{0.19\textwidth}>{\raggedright\arraybackslash}p{0.28\textwidth}>{\raggedright\arraybackslash}p{0.45\textwidth}@{}}
    \toprule
    \textbf{Proposition} & \textbf{Finding} & \textbf{Key evidence} \\
    \midrule
    P1 Low awareness & Supported & 78.2\% low familiarity; 22.2\% exposure \\
    \addlinespace
    P2 Value despite low awareness & Consistent; not falsifiable by this design & Low-familiarity majority still values highly (IT 4.21, security 4.08) \\
    \addlinespace
    P3 Variation by field and experience & Supported on both & Field $d=1.64$--$1.94$ (95\% CI above 1.4); mid-career experience-group effect $b=-0.40$/$-0.53$, $p<.001$ \\
    \addlinespace
    P4 Familiarity linked to value & Supported & $\rho=0.63$/$0.62$; standardized $\beta=0.34$/$0.33$ \\
    \bottomrule
  \end{tabular}
  \smallskip

  {\footnotesize \textit{Note.} All effect sizes are from the authenticated analytic sample ($N=1{,}411$); the field-gap intervals are 95\% bootstrap intervals (Table~\ref{tab:robustness_extra}). Propositions are supported in the sense of consistency with the cross-sectional evidence rather than causal confirmation.}
\end{table}

Across the three research questions the data tell a coherent story with three important qualifications. For Research Question~1, self-reported familiarity and recent exposure are low (familiarity mean near 1.9; 22.2\% exposure), yet endorsement of specific capability descriptions is near-universal (93.9\% agreement on Section~2). This gap is best characterized as low awareness coexisting with high capability-endorsement, rather than evidence of awareness as such.

For Research Questions~2 and~3, perceived value and anticipated benefit are high across the applicability, IT, and security blocks. The confidence intervals are narrow conditional on the accessible sample rather than on a random draw of the profession, and the two investment-willingness items are at ceiling, supplying conviction but not discrimination.

A very large field gap favoring Information Technology over Information Assurance dominates the comparative analyses, together with a non-monotonic experience pattern in which mid-career (11--15 year) respondents are the most reserved. The focused breakout analysis shows these two findings are largely independent rather than confounded: professional field and experience are only weakly associated, so the mid-career trough is a genuine experience-group difference present within both focal domains rather than a field-composition artifact. The Information Technology--Information Assurance gap likewise persists within each of the three interpretable experience cohorts (1--5, 6--10, and 11--15 years), which indicates the field difference is not merely an experience effect in disguise. The 16+ cohort, at $n=16$ across the two focal groups, is too small to interpret but shows Information Technology higher on all three composites as well.

Finally, the bias diagnostics are reassuring: the priming diagnostic argues against a pure question-framing explanation for the high endorsement, and the wave analysis returns a null result, with negligible differences between early- and late-arriving respondents on all three composites.

Read against the four propositions advanced in Chapter~1, the data support all four, with one refinement. P1 (low awareness) is clearly sustained: familiarity and recent exposure are both low. P2 (value despite low awareness) receives affirmative but qualified support, because capability endorsement and perceived value stay high even where familiarity is absent. This second proposition is, however, the most weakly tested of the four, since agreement with positively framed descriptions of capabilities tied to problems respondents already face is difficult to disconfirm by design. Its confirmation is therefore better read as consistent with the proposition than as a decisive test of it.

Stronger support for the study's contribution comes from three findings that could have failed and did not. In the priming diagnostic, high-familiarity and low-familiarity respondents diverged sharply rather than converging, which is the opposite of what a pure question-framing account predicts. High perceived value then persisted within the low-familiarity majority. The non-monotonic mid-career reservation, finally, was unanticipated in advance and stayed stable once the models adjusted for professional field.

Instead, the study treats the large professional-field gap as its most provisional result: it is stable across every statistical control the design permits, but the recruitment-frame difference that could confound it cannot be tested in these data (Table~\ref{tab:threats}). The chapter therefore presents it as a target for replication rather than as settled support for the contribution. P3 (variation by field and experience) finds support on both axes. The variation by field is large and stable across the applied statistical controls, and persists within each of the three interpretable experience cohorts, while the variation by experience is a genuine mid-career effect that survives adjustment for field. P4 (familiarity associated with higher value) rounds out the set: familiarity positively predicts perceived value without displacing the high baseline endorsement even unfamiliar respondents express.

Chapter~5 interprets these findings against the literature reviewed in Chapter~2, draws conclusions for each research question, addresses the limitations the diagnostics surface, and proposes directions for future research.
